\documentclass[11pt]{article}
\usepackage[margin=1in]{geometry}
\usepackage{amsmath,amssymb,amsthm,booktabs,hyperref}
\hypersetup{hidelinks}
\renewcommand{\labelitemi}{$\bullet$}
\newtheorem{theorem}{Theorem}
\newtheorem{proposition}[theorem]{Proposition}
\newtheorem{lemma}[theorem]{Lemma}
\newtheorem{corollary}[theorem]{Corollary}
\newtheorem{lemmaL}{Lemma}
\renewcommand{\thelemmaL}{\Alph{lemmaL}}
\newtheorem{propL}[lemmaL]{Proposition}
\theoremstyle{remark}
\newtheorem{remark}[theorem]{Remark}
\newcommand{\E}{\mathbb E}
\newcommand{\eqd}{\overset{d}{=}}
\title{Sharp $L_p$--$L_2$ moment comparison for centred log-concave random variables, with equality cases}
\author{Anonymous}
\date{15 October 2026}
\begin{document}
\maketitle

\begin{abstract}
For a nondegenerate mean-zero log-concave random variable $X$, we determine the sharp constants in $\|X\|_p\ge m_p\|X\|_2$ ($1<p<2$) and $\|X\|_p\le M_p\|X\|_2$ ($p>2$). The extremiser changes from the symmetric Laplace law to a centred exponential law at $p^*=1.6824572\ldots$, the unique solution in $(1,2)$ of $\Gamma(p+1)^{1/p}/\sqrt2=(\E|E-1|^p)^{1/p}$, with $E$ standard exponential. This proves the variance-normalised moment conjecture of Melbourne, Roysdon, Tang and Tkocz. Away from $p^*$ and $2$, the extremiser is unique up to nonzero scaling; at $p^*$ there are exactly the two stated families. At $p=2$ every admissible law is extremal. We give a self-contained, equality-sensitive proof of Eitan's real-exponent reduction using crossing arguments and an explicit Chebyshev system. The remaining two-variable optimisation is proved analytically for $p\ge7$ and by rigorous ball-arithmetic enclosures for $1<p<7$, with separate derivative estimates at the degenerate boundaries.
\end{abstract}

\section{Introduction}
For a random variable $X$ write $\|X\|_p=(\E|X|^p)^{1/p}$. Let $\mathcal L$ be the class of nondegenerate mean-zero log-concave random variables, that is, those with a density $e^{-V}$ for a proper convex $V:\mathbb R\to(-\infty,+\infty]$ (so supports may be half-lines or intervals; $E-1\in\mathcal L$). Melbourne, Roysdon, Tang and Tkocz \cite{MRTT} determined the sharp constants comparing $\|X\|_p$ with $\|X\|_1$ on $\mathcal L$, finding a phase transition in the extremiser. In their concluding remarks \cite[Sec.~3.1]{MRTT} they ask for the sharp constants in the variance-normalised problem,
\[
m_p=\inf_{X\in\mathcal L}\frac{\|X\|_p}{\|X\|_2}\quad(1<p<2),\qquad
M_p=\sup_{X\in\mathcal L}\frac{\|X\|_p}{\|X\|_2}\quad(p>2),
\]
and, based on numerical experiments, conjecture that the extremiser is the symmetric exponential for $1<p<p^*\approx1.68$ and the one-sided exponential for $p>p^*$. The phase transition and the value $p^*\approx1.68$ were thus found numerically by Melbourne, Roysdon, Tang and Tkocz; our contribution is a proof. Eitan \cite[Theorem~1.4]{Eitan} (Theorem~4 in the arXiv version) proved the case of even integers $p>2$.

At $p=0$ we use the limiting convention $\|X\|_0=\exp(\E\log|X|)$. The logarithm is integrable for the laws considered here: the density is bounded near zero and has exponential tails. The new results below concern $p>1$.

Let $E,E'$ be independent standard exponential random variables, $L=E-E'$ (symmetric exponential) and $F(p)=\E|E-1|^p$. Then $\|L\|_p/\|L\|_2=\Gamma(p+1)^{1/p}/\sqrt2$ and $\|E-1\|_p/\|E-1\|_2=F(p)^{1/p}$.

\begin{theorem}\label{thm:main}
There is a unique $p^*\in(1,2)$ with $\Gamma(p^*+1)^{1/p^*}/\sqrt2=F(p^*)^{1/p^*}$; numerically $p^*=1.6824572084\ldots$. Moreover:
\begin{itemize}
\item for $1<p\le p^*$, $m_p=\Gamma(p+1)^{1/p}/\sqrt2$, attained by the symmetric exponential;
\item for $p^*\le p<2$, $m_p=F(p)^{1/p}$, attained by the one-sided exponential $E-1$;
\item for $p>2$, $M_p=F(p)^{1/p}$, attained by $E-1$.
\end{itemize}
\end{theorem}

With \cite[Theorem~1]{MRTT}, which treats $-1<p\le1$ (symmetric exponential extremal), Theorem~\ref{thm:main} gives the sharp constants comparing $L_p$ with $L_2$ on $\mathcal L$ for all $p>-1$; for even $p>2$ it recovers \cite[Theorem~1.4]{Eitan}. We also classify the cases of equality. Here $X\eqd Y$ means equality in law.

\begin{theorem}[Equality cases]\label{thm:U}
Let $X\in\mathcal L$, and let $p^*$, $m_p$, $M_p$ be as in Theorem~\ref{thm:main}.
\begin{enumerate}
\item[(i)] If $1<p<p^*$, then $\|X\|_p=m_p\|X\|_2$ if and only if $X\eqd\lambda(E-E')$ for some $\lambda>0$.
\item[(ii)] If $p=p^*$, then $\|X\|_p=m_p\|X\|_2$ if and only if $X\eqd\lambda(E-E')$ or $X\eqd\pm\lambda(E-1)$ for some $\lambda>0$.
\item[(iii)] If $p^*<p<2$, then $\|X\|_p=m_p\|X\|_2$ if and only if $X\eqd\pm\lambda(E-1)$ for some $\lambda>0$.
\item[(iv)] If $p>2$, then $\|X\|_p=M_p\|X\|_2$ if and only if $X\eqd\pm\lambda(E-1)$ for some $\lambda>0$.
\end{enumerate}
\end{theorem}

\paragraph{Contribution and proof architecture.} The new optimisation and complete classification in Theorems~\ref{thm:main}--\ref{thm:U} concern the remaining real exponents. Eitan's reduction already covers real $1<a<b$; neither that scope nor the crossing method is claimed as new. Section~\ref{sec:selfred} supplies a self-contained, equality-sensitive proof of the reduction, following Eitan \cite[Sections~2--3]{Eitan} and the crossing methods of Eskenazis, Nayar and Tkocz \cite{ENT}. Sections~\ref{sec:largep}--\ref{sec:assembly} establish the reduced optimisation, and Section~\ref{sec:equality} extracts the strictness needed for the equality classification. All arguments and proof-critical computational material are included in this paper and its accompanying public package; no earlier draft is required.

\paragraph{Related work.} Murawski \cite{Murawski} proved the mean-zero bound $\|X\|_p\le\frac pq\|X\|_q$ ($p\ge q\ge1$) and sharp shifted-exponential results for general (not necessarily centred) log-concave variables. Eitan conjectured \cite[Conjecture~16 in arXiv v1]{Eitan} that for even $p<q$ the extremal $Z_s$ is $s\in\{0,1\}$; Theorem~\ref{thm:main} shows that this endpoint behaviour fails for real exponents, since for $(p,2)$ with $p<p^*$ the extremal parameter is $s=\frac12$. Theorem~\ref{thm:R2} reduces that conjecture, with its equality case, to an inequality within the family $Z_s$; we do not address it. Since $Z_s/\|Z_s\|_2$ is a two-term centred sum of exponentials, the problem also belongs to the study of sharp moment bounds for sums of exponential random variables \cite{BTT,BP}.

\paragraph{Method.} By Theorem~\ref{thm:R2} (Section~\ref{sec:selfred}) the extremal ratio is attained within the family $Z_s=s(E-1)-(1-s)(E'-1)$, $0\le s\le1$, and only there, so the theorem is equivalent to a two-parameter inequality for $G(s,p)=\frac1p\log\E|Z_s|^p-\frac12\log\E Z_s^2$ (Section~\ref{sec:reduction}). The inequality is degenerate along $s=0$ and $s=1/2$ (where $\partial_sG$ vanishes), along $p=2$ (where $G\equiv0$), and at the crossing $p=p^*$. We prove it analytically for $p\ge7$ (Section~\ref{sec:largep}). For $1<p<7$ we partition $[0,\frac12]\times(1,7]$ into regions with explicit margins, certified by ball arithmetic, and regions around the degenerate lines, where second-order lemmas reduce the claim to certified sign conditions on derivatives (Sections~\ref{sec:local}--\ref{sec:bulk}). All the certified inequalities are strict, and this gives the equality cases (Section~\ref{sec:equality}).

\paragraph{A concrete sharp bound.} At $p=4$, direct integration gives
$F(4)=\E(E-1)^4=24-24+12-4+1=9$. Thus every
$X\in\mathcal L$ satisfies $\E X^4\le9(\E X^2)^2$, with equality
exactly for a scaled or reflected centred exponential. This even-integer
case was already known from Eitan; it illustrates the normalisation rather
than a new fourth-moment inequality. In contrast, at $p=3/2<p^*$ the
sharp lower ratio is $\Gamma(5/2)^{2/3}/\sqrt2$ and only the Laplace
family attains it. At $p=2$ the ratio equals one for every admissible law.
Section~\ref{sec:stability} also shows that a nearly extremal ratio at a
fixed $p>1$, $p\ne2$, forces closeness in total variation and every finite
Wasserstein distance to the corresponding variance-normalised equality set.
This qualitative consequence has no explicit deficit-to-distance rate.

\section{A self-contained reduction (after Eitan)}\label{sec:selfred}
Throughout, $E,E'$ are independent standard exponential random variables and
\[
Z_s=s(E-1)-(1-s)(E'-1),\qquad 0\le s\le1.
\]
Thus $Z_{1/2}=\frac12(E-E')$, $Z_1=E-1$, $Z_0=1-E'\eqd-(E-1)$, and $Z_{1-s}\eqd-Z_s$ (exchange $E$ and $E'$).

\paragraph{Relation to Eitan's work.} The two theorems below give self-contained versions of Theorems~12 and~14 of Eitan's arXiv:2110.14382v1 \cite{Eitan}; Theorem~12 corresponds to Theorem~3.1 of the published article. The real-exponent reduction is already established there. Although its arXiv Theorem~12 does not state equality explicitly, Theorem~14 does, and the proof of Theorem~4 uses an equality characterisation for the reduction. Thus we do not claim priority for either the reduction or equality-sensitive reasoning. Our detailed comparison uses the arXiv text; the published full text was not inspected. The proofs below follow Eitan's crossing method and its antecedents \cite{ENT}, using Rolle's theorem, determinants, convexity, dominated convergence and the intermediate value theorem. Their role is to make the logical route to the complete extremiser classification explicit.

\begin{theorem}[Reduction with equality]\label{thm:R2}
Let $1<a<b$ be real and $X\in\mathcal L$. There exist $s_0\in[0,1]$ and $\lambda>0$ such that $\E|X|^a=\E|\lambda Z_{s_0}|^a$ and
\[
\E|X|^b\le\E|\lambda Z_{s_0}|^b,
\]
with equality if and only if $X\eqd\lambda Z_{s_0}$. Consequently the maximum below exists and
\[
\frac{\|X\|_b}{\|X\|_a}\le\max_{0\le s\le1}\frac{\|Z_s\|_b}{\|Z_s\|_a},
\]
with equality if and only if $X\eqd\lambda Z_{s_1}$ for some $\lambda>0$ and some $s_1$ at which the maximum is attained.
\end{theorem}

\begin{theorem}[Odd moments]\label{thm:R1}
Let $1<a<r$ be real and $X\in\mathcal L$. Then
\[
\frac{\E[|X|^r\operatorname{sgn}X]}{\|X\|_a^{\,r}}\le\frac{\E[|E-1|^r\operatorname{sgn}(E-1)]}{\|E-1\|_a^{\,r}},
\]
with equality if and only if $X\eqd\lambda(E-1)$ for some $\lambda>0$.
\end{theorem}

\subsection{Tools}

\begin{lemmaL}[log-concave preliminaries]\label{lem:A}
Let $X\in\mathcal L$ have density $e^{-V}$, where $V:\mathbb R\to(-\infty,\infty]$ is convex. Then $I:=\{V<\infty\}$ is an interval with nonempty interior; $V$ is finite and convex on $\operatorname{int}I$, and $f:=e^{-V}\mathbf 1_{\operatorname{int}I}$ is a version of the density; and there are $C,\kappa>0$ with $f(x)\le Ce^{-\kappa|x|}$ for all $x$. In particular $\E|X|^t<\infty$ for every $t\ge0$.
\end{lemmaL}
\begin{proof}
$I$ is convex and has positive Lebesgue measure because $\int e^{-V}=1$, so it is an interval with nonempty interior; its boundary has at most two points, a null set. Fix $x_0\in\operatorname{int}I$ and $m\in\partial V(x_0)$ (the subdifferential is nonempty at interior points). Then $V(x)\ge V(x_0)+m(x-x_0)$ for all $x\in\mathbb R$, so $f(x)\le e^{-V(x_0)}e^{-m(x-x_0)}$ everywhere; in particular $f$ is bounded on every bounded set, and on a bounded side of $I$ it vanishes beyond the endpoint.

Suppose $I$ is unbounded above. If $V(x)\le V(x_0)$ for all $x>x_0$, then $e^{-V}\ge e^{-V(x_0)}$ on $(x_0,\infty)$, contradicting integrability. So there is $x_1>x_0$ with $V(x_1)>V(x_0)$, and by convexity, for $x\ge x_1$,
\[
V(x)\ge V(x_1)+\kappa_+(x-x_1),\qquad \kappa_+=\frac{V(x_1)-V(x_0)}{x_1-x_0}>0.
\]
The same argument applies below if $I$ is unbounded below. Hence on each unbounded side $f$ decays exponentially beyond some point; on each bounded side $f$ vanishes beyond the endpoint; and in between $f$ is bounded by the affine-minorant bound. Combining these gives $C,\kappa>0$ with $f(x)\le Ce^{-\kappa|x|}$.
\end{proof}

\begin{lemmaL}[Descartes' rule for generalised polynomials]\label{lem:B}
Let $\lambda_0<\dots<\lambda_n$ be real and $c=(c_0,\dots,c_n)\neq0$. Then $P(t)=\sum_ic_it^{\lambda_i}$ has at most $\mathrm V(c)$ distinct zeros in $(0,\infty)$, where $\mathrm V(c)$ is the number of sign changes in $c$ after zero entries are deleted.
\end{lemmaL}
\begin{proof}
Induction on $\mathrm V(c)$. If $\mathrm V(c)=0$, all nonzero terms $c_it^{\lambda_i}$ have the same sign for $t>0$, and at least one is nonzero, so $P(t)\ne0$. If $\mathrm V(c)\ge1$, choose consecutive nonzero entries $c_j,c_{j'}$ ($j<j'$) of opposite signs and $\mu\in(\lambda_j,\lambda_{j'})$. The function $\psi(t)=t^{-\mu}P(t)$ has the same zeros as $P$ on $(0,\infty)$, and
\[
\psi'(t)=\sum_ic_i(\lambda_i-\mu)t^{\lambda_i-\mu-1}.
\]
Its exponents are distinct, and its coefficients are nonzero exactly where $c_i\ne0$. Passing from $c$ to these coefficients flips the sign of every entry with $i\le j$ and keeps every entry with $i\ge j'$; this preserves the sign changes inside each group and removes the one between $c_j$ and $c_{j'}$. So the coefficients of $\psi'$ have $\mathrm V(c)-1$ sign changes. If $P$ had $N$ distinct zeros, Rolle's theorem would give $\psi'$ at least $N-1$ distinct zeros, and by induction $N-1\le\mathrm V(c)-1$.
\end{proof}
This is the classical rule of Laguerre \cite[Part~V, Ch.~1, Sec.~6]{PolyaSzego}; the proof is included for completeness.

\begin{lemmaL}\label{lem:C}
For $c\in\mathbb R^{n+1}\setminus\{0\}$ let $c^*_i=(-1)^ic_i$. Then $\mathrm V(c)+\mathrm V(c^*)\le n-\mathbf 1[c_0=0]$.
\end{lemmaL}
\begin{proof}
Consider consecutive nonzero entries at indices $i<i'$, with gap $g=i'-i\ge1$. This pair contributes $\mathbf 1[c_ic_{i'}<0]+\mathbf 1[(-1)^gc_ic_{i'}<0]$ to the left side. If $g$ is odd the contribution is exactly $1\le g$; if $g$ is even it is $0$ or $2\le g$. Summing over pairs gives at most $i_{\rm last}-i_{\rm first}\le n-\mathbf 1[c_0=0]$.
\end{proof}

\begin{propL}[Chebyshev property]\label{prop:D}
Let $1<a<r<b$ and set
\[
u_0=1,\quad u_1=x,\quad u_2=|x|^a,\quad u_3=|x|^r\operatorname{sgn}x,\quad u_4=|x|^b .
\]
These are continuous on $\mathbb R$. For each $n\in\{0,\dots,4\}$, every nontrivial combination $\varphi=\sum_{i\le n}c_iu_i$ has at most $n$ distinct zeros in $\mathbb R$.
\end{propL}
\begin{proof}
Let $(\lambda_0,\dots,\lambda_4)=(0,1,a,r,b)$. For $t>0$, $u_i(t)=t^{\lambda_i}$ and $u_i(-t)=(-1)^it^{\lambda_i}$. So on $(0,\infty)$, $\varphi=\sum c_it^{\lambda_i}$ has at most $\mathrm V(c)$ distinct zeros by Lemma~\ref{lem:B}; $t\mapsto\varphi(-t)=\sum c_i^*t^{\lambda_i}$ has at most $\mathrm V(c^*)$ distinct zeros in $(0,\infty)$; and $0$ is a zero if and only if $c_0=0$. Lemma~\ref{lem:C} bounds the total by $n$.
\end{proof}
The same proof applies to any number of functions with alternating parity and increasing exponents $0<1<\lambda_2<\cdots$.

\begin{lemmaL}[alternating kernels]\label{lem:E}
Let $u_0,\dots,u_n$ be continuous on $\mathbb R$, and suppose every nontrivial combination of them has at most $n$ distinct zeros. Given $z_1<\dots<z_n$, there is a nontrivial combination $\varphi$ of $u_0,\dots,u_n$ with $\{\varphi=0\}=\{z_1,\dots,z_n\}$ that has constant nonzero sign on each interval $(z_j,z_{j+1})$ ($z_0=-\infty$, $z_{n+1}=+\infty$), these signs alternating in $j$.
\end{lemmaL}
\begin{proof}
For $t_0<\dots<t_n$ let $D(t_0,\dots,t_n)=\det[u_i(t_j)]_{i,j=0}^n$. It never vanishes: otherwise some $c\ne0$ satisfies $\sum_ic_iu_i(t_j)=0$ for all $j$, giving a nontrivial combination with $n+1$ zeros. Being continuous on the connected set $\{t_0<\dots<t_n\}$, $D$ has a constant sign $\sigma$. Put $\varphi(x)=\det[u_i(z_1),\dots,u_i(z_n),u_i(x)]$, a combination of the $u_i$ by expansion along the last column. It vanishes at each $z_j$. For $x\in(z_j,z_{j+1})$, moving the last column to position $j$ takes $n-j$ transpositions, so $\varphi(x)=(-1)^{n-j}D(z_1,\dots,z_j,x,z_{j+1},\dots,z_n)$, whose sign is $(-1)^{n-j}\sigma\ne0$. In particular $\varphi$ is nontrivial.
\end{proof}

\begin{lemmaL}[sign patterns and moment constraints]\label{lem:F}
Let $1<a<r<b$, let $u_i$ and $\lambda_i$ be as in Proposition~\ref{prop:D}, and let $n\in\{3,4\}$. Let $h$ be measurable and not a.e.\ zero, with $\int|h|(1+|x|^{\lambda_n})<\infty$ and $\int hu_i=0$ for $i=0,\dots,n-1$. Suppose there are $-\infty=y_0\le y_1\le\dots\le y_n\le y_{n+1}=+\infty$ such that $(-1)^{n-j}h\ge0$ a.e.\ on $(y_j,y_{j+1})$ for $j=0,\dots,n$ (alternating signs, the last one $+$). Then $\int hu_n>0$.
\end{lemmaL}
\begin{proof}
\emph{Step 1: a minimal alternating pattern.} Let $J=\{j:h$ is not a.e.\ zero on $(y_j,y_{j+1})\}=\{j_1<\dots<j_m\}$, which is nonempty, and let $e_l=(-1)^{n-j_l}$. Group the indices of $J$ into maximal runs of constant $e_l$; if $j_l$ is the last index of a run that is not the final run, put a change point at $y_{j_l+1}$. These points $z_1<\dots<z_k$ are finite and strictly increasing: $(y_{j_l},y_{j_l+1})$ has positive measure, so $y_{j_l+1}>y_{j_l}\ge-\infty$, and $y_{j_l+1}\le y_{j_{l+1}}<y_{j_{l+1}+1}\le+\infty$. Moreover $k\le m-1\le n$. Let $\omega$ be the step function equal to $e_1$ on $(-\infty,z_1)$ that switches sign at each $z_i$. Every interval $(y_j,y_{j+1})$ with $j\in J$ lies between the two change points that bound its run, so $\omega h\ge0$ a.e., and $\omega h$ is not a.e.\ zero.

\emph{Step 2: $k<n$ is impossible.} If $k\le n-1$, then $u_0,\dots,u_k$ satisfy the hypothesis of Lemma~\ref{lem:E} by Proposition~\ref{prop:D}. Lemma~\ref{lem:E} gives $\varphi\in\operatorname{span}\{u_0,\dots,u_k\}$ with, after a choice of sign, $\varphi\omega>0$ off $\{z_i\}$. Then $\varphi h\ge0$ a.e.\ and is not a.e.\ zero, so $\int\varphi h>0$; but $\int\varphi h=0$ by the constraints.

\emph{Step 3: $k=n$.} Then $m=n+1$ and $J=\{0,\dots,n\}$, so the final sign is $e_m=+1$. Take $\varphi$ from Lemma~\ref{lem:E} for $u_0,\dots,u_n$ with sign chosen so that $\varphi>0$ on $(z_n,\infty)$. Then $\varphi$ has the sign of $\omega$ off $\{z_i\}$, so $\int\varphi h>0$. Writing $\varphi=\sum_{i\le n}c_iu_i$, the constraints give $\int\varphi h=c_n\int hu_n$, so $c_n\ne0$. As $x\to+\infty$, $\varphi(x)/x^{\lambda_n}\to c_n$, because $u_n(x)=x^{\lambda_n}$ for $x>0$ has the largest exponent; since $\varphi>0$ on $(z_n,\infty)$, $c_n>0$. Hence $\int hu_n>0$.
\end{proof}

\begin{lemmaL}[crossing pattern against $Z_s$]\label{lem:G}
Let $X\in\mathcal L$ have the density $f=e^{-V}\mathbf 1_{\operatorname{int}I}$ of Lemma~\ref{lem:A}. Let $\lambda>0$, $s\in[0,1]$, let $g$ be the density of $\lambda Z_s$, and let $h=g-f$.
\begin{enumerate}
\item There are $-\infty\le y_1\le y_2\le y_3\le y_4\le+\infty$ such that, a.e., $h$ is $\ge0$, $\le0$, $\ge0$, $\le0$, $\ge0$ on the five successive intervals; this is the pattern of Lemma~\ref{lem:F} with $n=4$.
\item If $s=1$, there are $y_1\le y_2\le y_3$ with $h\le0$, $\ge0$, $\le0$, $\ge0$ a.e.\ on the four successive intervals; this is the pattern of Lemma~\ref{lem:F} with $n=3$.
\end{enumerate}
\end{lemmaL}
\begin{proof}
Put $c=1-s$ and $d=1-2s$. For $0<s<1$, $Z_s=d+(sE-cE')$, and $sE-cE'$ has density $e^{-t/s}$ for $t>0$ and $e^{t/c}$ for $t<0$ (the tails are given in Section~\ref{sec:reduction}). So $g>0$ everywhere and $\log g$ is affine on $L=(-\infty,\lambda d]$ and on $R=[\lambda d,\infty)$. For $s=0$, $Z_0=1-E'$ and $g$ is positive and log-affine on $(-\infty,\lambda]$ and zero after; for $s=1$, $Z_1=E-1$ and $g$ is positive and log-affine on $[-\lambda,\infty)$ and zero before.

\emph{Key observation.} Let $P$ be an interval on which $g=e^{\ell}$ with $\ell$ affine. On $P\cap\operatorname{int}I$, $h<0$ if and only if $\ell+V<0$; as $\ell+V$ is convex there, this set is an interval $N_P$ (possibly empty, possibly all of $P\cap\operatorname{int}I$, possibly half-infinite). On $P\setminus\overline I$, $h=g>0$, and $\partial I$ is null. Hence, a.e.\ on $P$, $h<0$ exactly on the interval $N_P$ and $h\ge0$ on the rest of $P$, which consists of at most two intervals, one on each side of $N_P$. Where $g=0$, $h=-f\le0$. Whether the endpoints of the intervals below are included is immaterial, since all statements are a.e.

The patterns follow.
\begin{itemize}
\item $0<s<1$: let $(y_1,y_2)$ be the endpoints of $N_L$ and $(y_3,y_4)$ those of $N_R$, so $y_2\le\lambda d\le y_3$. If $N_L=\emptyset$ take $y_1=y_2=\lambda d$; if $N_R=\emptyset$ take $y_3=y_4=\lambda d$.
\item $s=0$: let $(y_1,y_2)$ be the endpoints of $N_{(-\infty,\lambda]}$, $y_3=\lambda$ and $y_4=+\infty$ (the fifth interval is empty); if $N_{(-\infty,\lambda]}=\emptyset$ take $y_1=y_2=\lambda$.
\item $s=1$, part 1: $y_1=-\infty$, $y_2=-\lambda$, and $(y_3,y_4)$ the endpoints of $N_{[-\lambda,\infty)}$; if it is empty take $y_3=y_4=-\lambda$.
\item $s=1$, part 2: $y_1=-\lambda$ and $(y_2,y_3)$ the endpoints of $N_{[-\lambda,\infty)}$; if it is empty take $y_2=y_3=-\lambda$.
\end{itemize}
In each case the stated sign holds a.e.\ on each interval.
\end{proof}

\subsection{\texorpdfstring{Proofs of Theorems~\ref{thm:R2} and~\ref{thm:R1}}{Proofs of the reduction theorems}}

\begin{proof}[Proof of Theorem~\ref{thm:R1}]
Let $Y=\lambda(E-1)=\lambda Z_1$ with $\lambda>0$ chosen so that $\E|Y|^a=\E|X|^a$, let $g$ be its density and $h=g-f$. By Lemma~\ref{lem:A} all moments are finite, so $\int|h|(1+|x|^r)<\infty$. We have $\int h=0$ (both are densities), $\int xh=0$ (both are centred) and $\int|x|^ah=0$ (choice of $\lambda$); these are $\int hu_i=0$ for $i=0,1,2$, with $u_i$ as in Proposition~\ref{prop:D} for any $b>r$. By Lemma~\ref{lem:G}(2), $h$ has the $n=3$ pattern. If $h$ is not a.e.\ zero, Lemma~\ref{lem:F} with $n=3$ gives $\int|x|^r\operatorname{sgn}(x)\,h>0$, that is, $\E[|Y|^r\operatorname{sgn}Y]>\E[|X|^r\operatorname{sgn}X]$. Dividing by $\|X\|_a^r=\|Y\|_a^r$ and using the scale invariance of the ratio gives strict inequality unless $X\eqd Y$.
\end{proof}

\begin{proof}[Proof of Theorem~\ref{thm:R2}]
Fix $r=(a+b)/2$ and let $\rho(Y)=\E[|Y|^r\operatorname{sgn}Y]/(\E|Y|^a)^{r/a}$. It is scale invariant and $\rho(-Y)=-\rho(Y)$. Theorem~\ref{thm:R1} applied to $X$ and to $-X$ (both in $\mathcal L$) gives $\rho(X)\le\rho(Z_1)$ and $-\rho(X)\le\rho(Z_1)$. Since $Z_0\eqd-Z_1$, $\rho(Z_0)=-\rho(Z_1)$, so
\[
\rho(Z_0)\le\rho(X)\le\rho(Z_1).
\]
The map $s\mapsto Z_s$ is pointwise continuous and $|Z_s|\le E+E'+2$. For each $t\ge0$, $(E+E'+2)^t$ is integrable, so by dominated convergence $s\mapsto\E|Z_s|^t$ and $s\mapsto\E[|Z_s|^r\operatorname{sgn}Z_s]$ are continuous on $[0,1]$; moreover $\E|Z_s|^t>0$. Hence $s\mapsto\rho(Z_s)$ is continuous, and the intermediate value theorem gives $s_0\in[0,1]$ with $\rho(Z_{s_0})=\rho(X)$.

Let $\lambda=\|X\|_a/\|Z_{s_0}\|_a$ and $Y=\lambda Z_{s_0}$, with density $g_Y$. Then $\E|Y|^a=\E|X|^a$ and $\E[|Y|^r\operatorname{sgn}Y]=\rho(X)(\E|X|^a)^{r/a}=\E[|X|^r\operatorname{sgn}X]$. With $h=g_Y-f$ we have $\int|h|(1+|x|^b)<\infty$ and $\int hu_i=0$ for $i=0,\dots,3$. By Lemma~\ref{lem:G}(1), $h$ has the $n=4$ pattern, and Lemma~\ref{lem:F} gives $\E|Y|^b-\E|X|^b=\int|x|^bh>0$ unless $h=0$ a.e., that is, unless $X\eqd Y$.

\emph{Consequences.} Let $Q(s)=\|Z_s\|_b/\|Z_s\|_a$. By the continuity just proved (with $t=a$ and $t=b$), $Q$ is continuous on the compact interval $[0,1]$ and attains its maximum. Then
\[
\frac{\|X\|_b}{\|X\|_a}\le\frac{\|Y\|_b}{\|Y\|_a}=Q(s_0)\le\max Q.
\]
If $\|X\|_b/\|X\|_a=\max Q$, both inequalities are equalities: the first forces $X\eqd\lambda Z_{s_0}$, and the second shows that $s_0$ is a maximiser. Conversely, if $X\eqd\lambda Z_{s_1}$ with $s_1$ a maximiser, equality is immediate.
\end{proof}

\begin{remark}
If $s_0\in\{0,1\}$ in the proof of Theorem~\ref{thm:R2}, then $h$ has at most three sign changes, and Lemma~\ref{lem:F} forces $h=0$ a.e. For $s_0=1$, the pattern of Lemma~\ref{lem:G}(2) ($-,+,-,+$) is the $n=4$ pattern of Lemma~\ref{lem:F} with an empty first interval, and Step~2 of the proof of Lemma~\ref{lem:F} ($k\le3<4$) rules out $h\ne0$. For $s_0=0$ the pattern is $+,-,+,-$, so Lemma~\ref{lem:F} is applied in the same way to $-h$, which satisfies the same four constraints. This is consistent with the equality case of Theorem~\ref{thm:R1}.
\end{remark}

\section{Reduction to a two-parameter inequality}\label{sec:reduction}
Each $Z_s$ is itself in $\mathcal L$ (for $0<s<1$ it has the density of an asymmetric Laplace law shifted to mean zero), so the maximum in Theorem~\ref{thm:R2} is the sharp constant. Applying Theorem~\ref{thm:R2} with $(a,b)=(p,2)$ for $1<p<2$ and $(2,p)$ for $p>2$, and using $Z_{1-s}\overset d=-Z_s$:
\[
\log m_p=\min_{0\le s\le1/2}G(s,p)\quad(1<p<2),\qquad \log M_p=\max_{0\le s\le1/2}G(s,p)\quad(p>2),
\]
where $G(s,p)=\frac1p\log M(s,p)-\frac12\log N(s)$, $M(s,p)=\E|Z_s|^p$ and $N(s)=\E Z_s^2=s^2+(1-s)^2$. Since $G(0,p)=\frac1p\log F(p)=:G_0(p)$ and $G(\frac12,p)=\frac1p\log\Gamma(p+1)-\frac12\log2=:G_h(p)$, Theorem~\ref{thm:main} is equivalent to:
\begin{itemize}
\item[(T$_A$)] $G(s,p)\ge G_h(p)$ for $1<p\le p^*$, $0\le s\le\frac12$;
\item[(T$_B$)] $G(s,p)\ge G_0(p)$ for $p^*\le p<2$, $0\le s\le\frac12$;
\item[(T$_C$)] $G(s,p)\le G_0(p)$ for $p>2$, $0\le s\le\frac12$;
\end{itemize}
together with the existence and uniqueness of $p^*$.

\paragraph{A representation.} Put $c=1-s$, $d=1-2s$, and let $R$ be standard exponential. For $0<s\le\frac12$, $sE-cE'$ has tails $\mathbb P(sE-cE'>t)=se^{-t/s}$ and $\mathbb P(sE-cE'<-t)=ce^{-t/c}$ ($t\ge0$), the same as the mixture equal to $sR$ with probability $s$ and $-cR$ with probability $c$. Since $Z_s=sE-cE'+d$,
\begin{equation}\label{eq:rep}
M(s,p)=s\,\E(d+sR)^p+c\,\E|d-cR|^p,
\end{equation}
and $\E|x-R|^p=e^{-x}\bigl(\Gamma(p+1)+\int_0^xu^pe^u\,du\bigr)$ for $x\ge0$. In particular $F(p)=e^{-1}\bigl(\Gamma(p+1)+\int_0^1u^pe^u\,du\bigr)$.

\section{Large exponents}\label{sec:largep}
\begin{proposition}\label{prop:largep}
For every real $p\ge7$ and $0<s\le\frac12$,
\[
\begin{gathered}
N(s)^{p/2}F(p)-M(s,p)\;\ge\;\frac{s^2(1-s)^{p-1}}{e}\,K(p)>0,\\
K(p)=\Bigl(\frac p2-2e+2\Bigr)\Gamma(p+1)-\frac{e(p-2)}2-\frac{e(e-2)}{p+1}.
\end{gathered}
\]
In particular $G(s,p)<G_0(p)$, so (T$_C$) holds for $p\ge7$.
\end{proposition}
\begin{proof}
Write $\Gamma_p=\Gamma(p+1)$ and $J(x)=\int_0^xu^pe^u\,du$. Put $\tau=s/c\in(0,1]$, so that $c=1/(1+\tau)$, $d/c=1-\tau$ and $N(s)=c^2(1+\tau^2)$. By \eqref{eq:rep},
\[
\begin{aligned}
M(s,p)&=\frac{c^p}{1+\tau}\Bigl[\tau\,\E(1-\tau+\tau R)^p+e^{\tau-1}\bigl(\Gamma_p+J(1-\tau)\bigr)\Bigr],\\
N(s)^{p/2}F(p)&=c^p(1+\tau^2)^{p/2}e^{-1}\bigl(\Gamma_p+J(1)\bigr).
\end{aligned}
\]
With $A(\tau)=(1+\tau)(1+\tau^2)^{p/2}$ this gives $e(1+\tau)c^{-p}\bigl[N^{p/2}F-M\bigr]=S(\tau)$, where
\[
S(\tau)=A(\tau)\bigl(\Gamma_p+J(1)\bigr)-e\tau\,\E(1-\tau+\tau R)^p-e^\tau\bigl(\Gamma_p+J(1-\tau)\bigr).
\]
Since $t\mapsto t^p$ is convex and $1-\tau+\tau R$ is a convex combination of $1$ and $R$, $\E(1-\tau+\tau R)^p\le(1-\tau)+\tau\Gamma_p$. Hence $S\ge\Delta$ with
\[
\Delta(\tau)=\bigl[A(\tau)-e^\tau-e\tau^2\bigr]\Gamma_p+\Bigl\{A(\tau)J(1)-e^\tau J(1-\tau)-e\tau(1-\tau)\Bigr\}.
\]
For $\tau\in[0,1]$ and $p\ge2$ we use:
\begin{itemize}
\item[(a)] $A(\tau)\ge(1+\tau)(1+\frac p2\tau^2)\ge1+\tau+\frac p2\tau^2$, by Bernoulli's inequality with exponent $p/2\ge1$;
\item[(b)] $0\le e^\tau-1-\tau\le(e-2)\tau^2$;
\item[(c)] $0\le J(1)\le e/(p+1)$;
\item[(d)] $\int_{1-\tau}^1u^pe^u\,du\ge e^{1-\tau}\bigl(1-(1-\tau)^{p+1}\bigr)/(p+1)$;
\item[(e)] $1-(1-\tau)^{p+1}\ge(p+1)\tau-\frac{(p+1)p}2\tau^2$, by Taylor's theorem, since the third derivative of $(1-\tau)^{p+1}$ is $\le0$ on $[0,1]$.
\end{itemize}
By (a) and (b), the first bracket is at least $(\frac p2-2e+2)\tau^2$. For the second, write $J(1-\tau)=J(1)-\int_{1-\tau}^1u^pe^u\,du$ and use (a), (b), (c), (d), (e):
\[
A(\tau)J(1)-e^\tau J(1-\tau)\ge(1+\tau-e^\tau)J(1)+e\,\frac{1-(1-\tau)^{p+1}}{p+1}\ge-\frac{e(e-2)}{p+1}\tau^2+e\Bigl(\tau-\frac p2\tau^2\Bigr).
\]
Subtracting $e\tau(1-\tau)$, the terms linear in $\tau$ cancel, and the second bracket is at least $-\bigl(\frac{e(e-2)}{p+1}+\frac{e(p-2)}2\bigr)\tau^2$. Hence $S(\tau)\ge\Delta(\tau)\ge\tau^2K(p)$, which is the stated bound after multiplying by $c^p/(e(1+\tau))$.

Finally let $p\ge7$. Then $\frac p2-2e+2>\frac{11}2-\frac{87}{16}=\frac1{16}$ (as $e<\frac{87}{32}$). Since $\Gamma$ is increasing on $[2,\infty)$, $\Gamma_p=p\,\Gamma(p)\ge720p$. Also $\frac{e(p-2)}2<\frac{87}{64}p$ and $\frac{e(e-2)}{p+1}<\frac{87}{32}\cdot\frac{23}{32}\cdot\frac18=\frac{2001}{8192}$. Thus
\[
K(p)>45p-\frac{87}{64}p-\frac{2001}{8192}=\frac{2793}{64}p-\frac{2001}{8192}>0. \qedhere
\]
\end{proof}
\begin{remark}
The argument needs $\frac p2>2e-2$, i.e.\ $p>4e-4\approx6.87$; indeed $K(6.8)<0<K(6.9)$. The loss comes from the chord bound for $\E(1-\tau+\tau R)^p$.
\end{remark}


\section{Derivatives in $s$ and the endpoints}\label{sec:local}
Let $W=E+E'-2=\partial_sZ_s$.
\begin{lemma}\label{lem:cond}
Let $0<s<1$ and $X=sE-cE'$, so that $Z_s=d+X$. Then $\mathbb P(X>0)=s$, and conditionally on $\{X>0\}$ the pair $(Z_s,W)$ has the law of $(d+sR,\,(R-1)+(T-1))$. Conditionally on $\{X<0\}$ it has the law of $(d-cR,\,(R-1)+(T-1))$. Here $R,T$ are independent standard exponential random variables. At $s=0,1$ use the direct formulas for $Z_s,W$, or their moment limits; no conditioning on a zero-probability event is intended.
\end{lemma}
\begin{proof}
On $\{sE>cE'\}$, condition on $E'$. By memorylessness $R:=E-cE'/s$ is standard exponential and independent of $E'$, while $E'$ has conditional density proportional to $e^{-u}e^{-cu/s}=e^{-u/s}$ at $u$, so $E'=sT$. Then $X=sR$ and $W=R+cT+sT-2$. The other side is symmetric.
\end{proof}

\begin{lemma}\label{lem:deriv}
For $p>1$, $G(\cdot,p)$ is $C^2$ on $[0,\frac12]$. With $\tau=-d/N$ and $V=W-\tau Z_s$, put $A=\E[|Z_s|^{p-2}Z_sV]/M$ and $B=\E[|Z_s|^{p-2}V^2]/M$. Then
\[
\partial_sG=A,\qquad \partial_s^2G=(p-1)B-pA^2+\frac{2d}{N}A-\frac1{N^2}.
\]
\end{lemma}
\begin{proof}
The map $s\mapsto Z_s$ is affine with slope $W$. Differentiation under the expectation is justified by dominated convergence for the first derivative. For the second, use the absolute continuity of $z\mapsto|z|^{p-2}z$, whose derivative $(p-1)|z|^{p-2}$ is locally integrable for $p>1$, together with Tonelli; finiteness comes from the closed forms in supplement \texttt{local.md}. This gives $M_s=p\E[|Z|^{p-2}ZW]$ and $M_{ss}=p(p-1)\E[|Z|^{p-2}W^2]$. Substituting $W=V+\tau Z$ and using $\tau'+\tau^2=1/N^2$ (as $d^2+4sc=1$) gives the formulas.
\end{proof}

\begin{proposition}\label{prop:endpoints}
For every $p>1$:
\begin{itemize}
\item[(i)] $\partial_sG(0,p)=0$ and $\partial_sG(\frac12,p)=0$;
\item[(ii)] $\partial_s^2G(0,p)=\frac{p-1}p\bigl(\frac1{F(p)}-1\bigr)$ and $\partial_s^2G(\frac12,p)=\frac{4(p-2)^2}p$.
\end{itemize}
\end{proposition}
\begin{proof}
At $s=0$, $c=d=1$ and $V=W+Z=E-1$ is independent of $Z=1-E'$, so $\E[|Z|^{p-2}ZV]=0$. Hence $A=0$ and $B=F(p-2)/F(p)$. Two integrations by parts against $e^{-t}$ give $p(p-1)F(p-2)=F(p)-1+p$, which yields the first formula in (ii).

At $s=\frac12$, $Z_{1-s}\overset d=-Z_s$ (exchange $E$ and $E'$) and $N(1-s)=N(s)$, so $G(\cdot,p)$ is even about $\frac12$, and $\partial_sG(\frac12,p)=0$. There $d=0$, $N=\frac12$, $M=2^{-p}\Gamma(p+1)$, $|Z|=R/2$ on both sides, and $\E[V\mid\cdot]=R-1$. So $(p-1)B=4(p-1)\bigl(\Gamma(p+1)-2\Gamma(p)+2\Gamma(p-1)\bigr)/\Gamma(p+1)=4(p-3+4/p)$, and $\partial_s^2G=4(p-3+4/p)-4=4(p-2)^2/p$.
\end{proof}

\begin{proposition}[computer-assisted]\label{prop:local}
\begin{itemize}
\item[(L1)] $\partial_s^2G>\frac{6}{125}$ on $[\frac9{20},\frac12]\times(1,\frac74]$.
\item[(L2)] $\partial_s^2G>\frac{29}{2000}$ on $[0,\frac1{20}]\times[\frac85,\frac{39}{20}]$.
\item[(L3)] $\partial_s^2G<-\frac{31}{2000}$ on $[0,\frac1{20}]\times[\frac{41}{20},7]$.
\end{itemize}
Consequently, $G\ge G_h$ on $[\frac9{20},\frac12]\times(1,\frac74]$, $G\ge G_0$ on $[0,\frac1{20}]\times[\frac85,\frac{39}{20}]$, and $G\le G_0$ on $[0,\frac1{20}]\times[\frac{41}{20},7]$.
\end{proposition}
\begin{proof}[Certificate]
By Lemma~\ref{lem:cond}, the quantities $M$, $\E[|Z|^{p-2}ZV]$ and $(p-1)\E[|Z|^{p-2}V^2]$ are explicit combinations of $\Gamma(p+1)$, $\Gamma(p)$, upper incomplete gamma functions, $\int_0^xu^ae^u\,du=\frac{x^{a+1}}{a+1}{}_1F_1(a+1;a+2;x)$ and ${}_1F_1(1;b;z)$ (supplement \texttt{local.md}, Lemma~4 and Proposition~B). The terms that are singular as $p\to1$ appear with a factor $p-1$ and are rewritten as $d^{p-1}\bigl[\frac N{c^2}{}_1F_1(1;p;-x)-\frac N{s^2}{}_1F_1(1;p;d/s)\bigr]$, which is continuous at the corner $(\frac12,1)$, where $\partial_s^2G=4$.

Each region is covered by boxes with exact rational endpoints. On each box, $p$ is a ball, the $s$-dependence is enclosed through monotone hulls of the special-function pieces, and $d^{p-1}$ is enclosed in $[d_0^{e_1},d_1^{e_0}]$. Precision is $128$ bits, boxes are bisected down to width $2^{-22}$, and an undecided comparison counts as a failure. Region L1 used $50{,}214$ boxes, L2 $129{,}875$ and L3 $57{,}496$, all without failure.

The consequences follow from Proposition~\ref{prop:endpoints}(i) and Taylor's theorem: $\partial_sG$ has a fixed sign on each strip away from the endpoint.
\end{proof}


\section{The neighbourhood of $p=2$}\label{sec:strip}
At $p=2$ we have $M(s,2)=N(s)$, so $G(\cdot,2)\equiv0$ and every comparison degenerates. Let $\mathcal H(s,p)=\partial_p\bigl[G(s,p)-G_0(p)\bigr]$.

\begin{proposition}[computer-assisted]\label{prop:strip}
\begin{itemize}
\item[(S1)] $\mathcal H<0$ on $[\frac1{20},\frac12]\times[\frac{39}{20},\frac{41}{20}]$.
\item[(S2)] $\partial_p\partial_s^2G<0$ on $[0,\frac1{20}]\times[\frac{39}{20},\frac{41}{20}]$.
\end{itemize}
Consequently, for $p\in[\frac{39}{20},\frac{41}{20}]$ and $0<s\le\frac12$, $G(s,p)-G_0(p)$ has the sign of $2-p$.
\end{proposition}
\begin{proof}
\emph{Consequence.} Since $G(0,\cdot)=G_0$, $\mathcal H(0,p)=0$. The mixed partial derivatives of $G$ up to $\partial_p\partial_s^2$ exist, are continuous and commute on the strip (supplement \texttt{strip.md}, Lemma~3), so $\partial_s\mathcal H(0,p)=\partial_p\partial_sG(0,p)=0$ by Proposition~\ref{prop:endpoints}(i), and $\partial_s^2\mathcal H=\partial_p\partial_s^2G<0$ on $[0,\frac1{20}]$ by (S2), so Taylor's formula gives $\mathcal H<0$ on $(0,\frac1{20}]$. Together with (S1), $\mathcal H<0$ for $0<s\le\frac12$. Since $G(s,2)-G_0(2)=0$, we have $G(s,p)-G_0(p)=\int_2^p\mathcal H(s,t)\,dt$, which has the sign of $2-p$.

\emph{Certificate.} Write $q=p-2$, $\Phi_A=\E|Z_s|^{2+q}$, $\Phi_B=\E[|Z_s|^qZ_s(R-1)]$ and $\Phi_C=\E[|Z_s|^q((R-1)^2+1)]$, the expectations being taken in the conditioned form of Lemma~\ref{lem:cond}. Then $M_s=p\Phi_B$, $M_{ss}=p(p-1)\Phi_C$, $\partial_s\partial_pG=\partial_q(\Phi_B/\Phi_A)$ and $\partial_p\partial_s^2G=\partial_q\bigl[(1+q)\Phi_C/\Phi_A-(2+q)(\Phi_B/\Phi_A)^2\bigr]$ (with the $N$-terms independent of $p$).

On the negative side ($Z\le d$) the moments have closed forms in $q$ via Gamma functions and a rapidly convergent series with bounded tail. On the positive side they are enclosed by order-$12$ Taylor polynomials in $q$, whose coefficients are rigorous integrals (Arb's \texttt{acb.integral} with explicit tail and endpoint bounds), plus an explicit remainder bound.

S1 is certified with a mean-value form in $s$ and $q$ ($157$ boxes, worst certified upper bound below $-1.5\times10^{-5}$). S2 is certified with direct evaluation in $s$ and a mean-value form in $q$ ($64$ boxes, worst bound below $-0.040$). The shipped outputs retain exact rational enclosure endpoints. Both certificates are regenerated in normal and optimised replay; the regression suite rejects altered signs, empty or overlapping coverage, and malformed domains.
\end{proof}


\section{The crossing point and the bulk}\label{sec:bulk}
Let $H(p)=G_0(p)-G_h(p)$, where $G_0(p)=\frac1p\log F(p)$ and $G_h(p)=\frac1p\log\Gamma(p+1)-\frac12\log2$.

\begin{proposition}[computer-assisted]\label{prop:crossing}
Let $p^*_{\rm lo}=1.682457208$ and $p^*_{\rm hi}=1.682457209$. Then:
\begin{itemize}
\item $H>0$ on $[1,p^*_{\rm lo}]$;
\item $H<0$ on $[p^*_{\rm hi},\frac{39}{20}]$;
\item $H'<-0.016$ on $[p^*_{\rm lo},p^*_{\rm hi}]$.
\end{itemize}
Hence $H$ has exactly one zero $p^*$ in $[1,\frac{39}{20}]$, $H>0$ on $[1,p^*)$, and $H<0$ on $(p^*,\frac{39}{20}]$.
\end{proposition}

\begin{proposition}[computer-assisted]\label{prop:bulk}
Let $\eta=\frac1{20}$. The following hold:
\begin{itemize}
\item[(B1)] $G(s,p)>G_h(p)$ for $s\in[0,\frac9{20}]$, $p\in[1,p^*_{\rm hi}-\eta]$;
\item[(B2)] $G(s,p)>\min\{G_0(p),G_h(p)\}$ for $s\in[\frac1{20},\frac9{20}]$, $p\in[p^*_{\rm lo}-\eta,p^*_{\rm hi}+\eta]$;
\item[(B3)] $G(s,p)>G_0(p)$ for $s\in[\frac1{20},\frac12]$, $p\in[p^*_{\rm lo}+\eta,\frac{39}{20}]$;
\item[(B4)] $G(s,p)<G_0(p)$ for $s\in[\frac1{20},\frac12]$, $p\in[\frac{41}{20},7]$.
\end{itemize}
\end{proposition}
\begin{proof}[Certificates]
$M(s,p)=s\,J_p+c^{p+1}e^{-x}\bigl(\Gamma(p+1)+\frac{x^{p+1}}{p+1}{}_1F_1(p+1;p+2;x)\bigr)$, where $x=d/c$ and $J_q=s^qe^{d/s}\Gamma(q+1,d/s)$. The partial derivatives $M_s$ and $M_p=\E|Z_s|^p\log|Z_s|$ have explicit forms (supplement \texttt{bulk.md}).

On a box with centre $(s_0,p_0)$ and half-widths $(h_s,h_p)$, the difference $D$ to be certified is enclosed by $D(s_0,p_0)+[\partial_pD]\,[-h_p,h_p]+[\partial_sD]\,[-h_s,h_s]$. Here $[\partial_pD]$ is bounded through the convexity of $p\mapsto\E|Z_s|^p$: $M_p=\E|Z_s|^p\log|Z_s|$ and $F'$ are nondecreasing in $p$ (their derivatives are $\E|Z_s|^p\log^2|Z_s|\ge0$), so their values at the exact rational endpoints bound them; $M$ and $F$ are enclosed separately and interval division is used. $[\partial_sD]$ is bounded from monotone enclosures of the closed-form pieces in $s$. No incomplete gamma function is evaluated at a wide parameter ball.

Branch and bound at $128$ bits gives:
\begin{itemize}
\item B1: $1653$ boxes, margin $\ge3.2\times10^{-6}$;
\item B2: $1127$ boxes, margin $\ge4.0\times10^{-8}$;
\item B3: $4413$ boxes, margin $\ge1.4\times10^{-8}$;
\item B4: $8635$ boxes, margin $\ge8.0\times10^{-8}$;
\item the crossing lemma: $73$, $65$ and $1$ boxes.
\end{itemize}
There were no failures. B1--B4 are regenerated at $128$ bits in both normal and optimised replay. The four deliberately false bulk statements are rejected in both modes. Optional floating-point quadrature checks are diagnostic only and do not replace the interval certificates.
\end{proof}


\section{\texorpdfstring{Proof of Theorem~\ref{thm:main}}{Proof of the main theorem}}\label{sec:assembly}
By Section~\ref{sec:reduction} it suffices to prove (T$_A$), (T$_B$), (T$_C$) and the uniqueness of $p^*$. Throughout, $s\in[0,\frac12]$.

\paragraph{The crossing.} By Proposition~\ref{prop:crossing}, $H=G_0-G_h$ has a unique zero $p^*$ in $[1,\frac{39}{20}]$, with $H>0$ before it and $H<0$ after it. On $[\frac{39}{20},2)$, Proposition~\ref{prop:strip} at $s=\frac12$ gives $G_h-G_0>0$, so $p^*$ is the unique zero in $(1,2)$.

\paragraph{(T$_A$): $1<p\le p^*$.}
\begin{itemize}
\item For $p\le p^*_{\rm hi}-\eta$: B1 covers $s\le\frac9{20}$ and L1 covers $s\ge\frac9{20}$.
\item For $p\in[p^*_{\rm hi}-\eta,p^*]$:
  \begin{itemize}
  \item for $s\le\frac1{20}$, L2 gives $G\ge G_0$, and $G_0\ge G_h$ because $p\le p^*$;
  \item for $s\in[\frac1{20},\frac9{20}]$, B2 gives $G>\min\{G_0,G_h\}=G_h$;
  \item for $s\ge\frac9{20}$, L1 applies.
  \end{itemize}
\end{itemize}

\paragraph{(T$_B$): $p^*\le p<2$.}
\begin{itemize}
\item For $p\in[p^*,p^*_{\rm lo}+\eta]$:
  \begin{itemize}
  \item L2 covers $s\le\frac1{20}$;
  \item B2 gives $G>\min\{G_0,G_h\}=G_0$ on $[\frac1{20},\frac9{20}]$;
  \item L1 gives $G\ge G_h\ge G_0$ on $[\frac9{20},\frac12]$.
  \end{itemize}
\item For $p\in[p^*_{\rm lo}+\eta,\frac{39}{20}]$: L2 covers $s\le\frac1{20}$ and B3 the rest.
\item For $p\in[\frac{39}{20},2)$: Proposition~\ref{prop:strip}.
\end{itemize}

\paragraph{(T$_C$): $p>2$.}
\begin{itemize}
\item $(2,\frac{41}{20}]$: Proposition~\ref{prop:strip}.
\item $[\frac{41}{20},7]$: L3 for $s\le\frac1{20}$ and B4 for the rest.
\item $[7,\infty)$: Proposition~\ref{prop:largep}. \qed
\end{itemize}


\section{\texorpdfstring{Equality cases: proof of Theorem~\ref{thm:U}}{Equality cases}}\label{sec:equality}

\paragraph{Reduction to $G$.} Recall $G(1-s,p)=G(s,p)$, since $Z_{1-s}\eqd-Z_s$ and $N(1-s)=N(s)$. For $1<p<2$ apply Theorem~\ref{thm:R2} with $(a,b)=(p,2)$, where $\|Z_s\|_2/\|Z_s\|_p=e^{-G(s,p)}$; for $p>2$ apply it with $(a,b)=(2,p)$, where $\|Z_s\|_p/\|Z_s\|_2=e^{G(s,p)}$. By the equality clause of Theorem~\ref{thm:R2}, equality in Theorem~\ref{thm:main} holds for $X\in\mathcal L$ if and only if $X\eqd\lambda Z_{s_1}$ for some $\lambda>0$ and some $s_1\in[0,1]$ that minimises $G(\cdot,p)$ on $[0,1]$ when $1<p<2$, or maximises it when $p>2$. Theorem~\ref{thm:U} therefore follows from Lemma~\ref{lem:S} below together with this symmetry. Lemma~\ref{lem:S} shows that the set of extremal parameters in $[0,1]$ is $\{\frac12\}$ in case (i), $\{0,\frac12,1\}$ in case (ii), and $\{0,1\}$ in cases (iii) and (iv). Finally, $\lambda Z_{1/2}=\frac\lambda2(E-E')$, $\lambda Z_1=\lambda(E-1)$ and $\lambda Z_0\eqd-\lambda(E-1)$.

\paragraph{Facts used.} Lemma~\ref{lem:S} uses only the following statements, all established above with strict inequalities on the stated closed sets. Here $\eta=\frac1{20}$, $p^*_{\rm lo}=1.682457208$, $p^*_{\rm hi}=1.682457209$, and $G_0$, $G_h$, $H=G_0-G_h$ are as in Sections~\ref{sec:reduction} and~\ref{sec:bulk}.
\begin{itemize}
\item[(F1)] Lemma~\ref{lem:deriv}: for each $p>1$, $G(\cdot,p)$ is $C^2$ on $[0,\frac12]$.
\item[(F2)] Proposition~\ref{prop:endpoints}(i): for each $p>1$, $\partial_sG(0,p)=\partial_sG(\frac12,p)=0$.
\item[(F3)] Proposition~\ref{prop:local}: (L1) $\partial_s^2G>\frac6{125}$ on $[\frac9{20},\frac12]\times(1,\frac74]$; (L2) $\partial_s^2G>\frac{29}{2000}$ on $[0,\frac1{20}]\times[\frac85,\frac{39}{20}]$; (L3) $\partial_s^2G<-\frac{31}{2000}$ on $[0,\frac1{20}]\times[\frac{41}{20},7]$.
\item[(F4)] Proposition~\ref{prop:strip}: for $p\in[\frac{39}{20},\frac{41}{20}]$ and $0<s\le\frac12$, $G(s,p)-G_0(p)=\int_2^p\mathcal H(s,t)\,dt$ with $\mathcal H<0$; hence $G(s,p)>G_0(p)$ for $\frac{39}{20}\le p<2$ and $G(s,p)<G_0(p)$ for $2<p\le\frac{41}{20}$.
\item[(F5)] Proposition~\ref{prop:crossing}: $H>0$ on $[1,p^*_{\rm lo}]$, $H<0$ on $[p^*_{\rm hi},\frac{39}{20}]$ and $H'<0$ on $[p^*_{\rm lo},p^*_{\rm hi}]$. Hence $p^*_{\rm lo}<p^*<p^*_{\rm hi}$, $H>0$ on $[1,p^*)$, $H(p^*)=0$ and $H<0$ on $(p^*,\frac{39}{20}]$.
\item[(F6)] Proposition~\ref{prop:bulk}, on closed rectangles: (B1) $G>G_h$ on $[0,\frac9{20}]\times[1,p^*_{\rm hi}-\eta]$; (B2) $G>\min\{G_0,G_h\}$ on $[\frac1{20},\frac9{20}]\times[p^*_{\rm lo}-\eta,p^*_{\rm hi}+\eta]$; (B3) $G>G_0$ on $[\frac1{20},\frac12]\times[p^*_{\rm lo}+\eta,\frac{39}{20}]$; (B4) $G<G_0$ on $[\frac1{20},\frac12]\times[\frac{41}{20},7]$. The certified minimum margins are $3.2\times10^{-6}$, $4.0\times10^{-8}$, $1.4\times10^{-8}$ and $8.0\times10^{-8}$.
\item[(F7)] Proposition~\ref{prop:largep}: for $p\ge7$ and $0<s\le\frac12$, $N(s)^{p/2}F(p)-M(s,p)>0$. Since $G(s,p)-G_0(p)=\frac1p\log\bigl(M(s,p)/(N(s)^{p/2}F(p))\bigr)$, this is $G(s,p)<G_0(p)$.
\end{itemize}

\setcounter{lemmaL}{18}
\begin{lemmaL}[strict separation]\label{lem:S}
For $s\in[0,\frac12]$:
\begin{itemize}
\item[(a)] if $1<p<p^*$, then $G(s,p)>G_h(p)$ for $s<\frac12$;
\item[(b)] if $p=p^*$, then $G(s,p^*)>G_0(p^*)=G_h(p^*)$ for $0<s<\frac12$;
\item[(c)] if $p^*<p<2$, then $G(s,p)>G_0(p)$ for $s>0$;
\item[(d)] if $p>2$, then $G(s,p)<G_0(p)$ for $s>0$.
\end{itemize}
\end{lemmaL}
\begin{proof}
\emph{Strict local consequences.} By (F1), (F2) and Taylor's formula with integral remainder, (F3) gives:
\begin{itemize}
\item on L1, for $\frac9{20}\le s<\frac12$ and $1<p\le\frac74$: $G(s,p)-G_h(p)=\int_s^{1/2}(t-s)\,\partial_s^2G(t,p)\,dt\ge\frac3{125}(\frac12-s)^2>0$;
\item on L2, for $0<s\le\frac1{20}$ and $\frac85\le p\le\frac{39}{20}$: $G(s,p)-G_0(p)=\int_0^s(s-t)\,\partial_s^2G(t,p)\,dt\ge\frac{29}{4000}s^2>0$;
\item on L3, for $0<s\le\frac1{20}$ and $\frac{41}{20}\le p\le7$: $G(s,p)-G_0(p)\le-\frac{31}{4000}s^2<0$.
\end{itemize}
Note $\frac85<p^*_{\rm hi}-\eta<p^*_{\rm lo}+\eta<\frac74<\frac{39}{20}$.

\emph{(a) $1<p<p^*$.} If $p\le p^*_{\rm hi}-\eta$, B1 gives $G>G_h$ for $s\le\frac9{20}$ and L1 for $\frac9{20}\le s<\frac12$. If $p\in[p^*_{\rm hi}-\eta,p^*)$: at $s=0$, $G=G_0>G_h$ by (F5); for $0<s\le\frac1{20}$, L2 gives $G>G_0>G_h$; for $s\in[\frac1{20},\frac9{20}]$, B2 gives $G>\min\{G_0,G_h\}=G_h$; for $s\in[\frac9{20},\frac12)$, L1 applies.

\emph{(b) $p=p^*$.} $G_0(p^*)=G_h(p^*)$ by (F5). For $0<s\le\frac1{20}$, L2 gives $G>G_0(p^*)$; for $s\in[\frac1{20},\frac9{20}]$, B2 gives $G>\min\{G_0,G_h\}(p^*)$; for $s\in[\frac9{20},\frac12)$, L1 gives $G>G_h(p^*)$.

\emph{(c) $p^*<p<2$.} If $p\in(p^*,p^*_{\rm lo}+\eta]$: for $0<s\le\frac1{20}$, L2; for $s\in[\frac1{20},\frac9{20}]$, B2 gives $G>\min\{G_0,G_h\}=G_0$ since $H<0$; for $s\in[\frac9{20},\frac12]$, L1 gives $G\ge G_h>G_0$. If $p\in[p^*_{\rm lo}+\eta,\frac{39}{20}]$: L2 for $0<s\le\frac1{20}$ and B3 for $s\ge\frac1{20}$. If $p\in[\frac{39}{20},2)$: (F4).

\emph{(d) $p>2$.} If $p\in(2,\frac{41}{20}]$: (F4). If $p\in[\frac{41}{20},7]$: L3 for $0<s\le\frac1{20}$ and B4 for $s\in[\frac1{20},\frac12]$. If $p\ge7$: (F7).

The region boundaries are those of the assembly in Section~\ref{sec:assembly}.
\end{proof}

\begin{remark}
The same bounds give quadratic lower bounds in $s$ near the extremal parameters; for example $G(s,p)-G_0(p)\ge\frac{29}{4000}s^2$ on $[0,\frac1{20}]\times[\frac85,\frac{39}{20}]$. Combined with Theorem~\ref{thm:R2}, this is a stability statement for the reduced parameter $s_0$ only. It is not stability in law, and it is not claimed as such.
\end{remark}


\section{Qualitative stability in distribution}\label{sec:stability}
The equality classification has the following compactness consequence.
The convergence results used here are standard; no quantitative
deficit-to-distance rate is asserted.

Let $\mathcal C$ be the laws of mean-zero, variance-one log-concave
random variables. Write
\[
 \nu_L=\operatorname{Law}((E-E')/\sqrt2),\qquad
 \nu_+=\operatorname{Law}(E-1),\qquad
 \nu_-=\operatorname{Law}(1-E),
\]
and define the variance-normalised equality set
\[
 \mathcal A_p=
 \begin{cases}
 \{\nu_L\},&1<p<p^*,\\
 \{\nu_L,\nu_+,\nu_-\},&p=p^*,\\
 \{\nu_+,\nu_-\},&p^*<p<2\text{ or }p>2.
 \end{cases}
\]
For $X$ with law $\mu\in\mathcal C$, put
\[
 \Delta_p(\mu)=
 \begin{cases}
 \log(\|X\|_p/m_p),&1<p<2,\\
 \log(M_p/\|X\|_p),&p>2.
 \end{cases}
\]
This is nonnegative by Theorem~\ref{thm:main}. For $1\le q<\infty$,
use the conventions
\[
 d_{\rm TV}(\mu,\nu)=\sup_{B\text{ Borel}}|\mu(B)-\nu(B)|,
 \qquad
 W_q(\mu,\nu)=\inf_{(X,Y)}(\E|X-Y|^q)^{1/q},
\]
where the infimum is over couplings with laws $\mu,\nu$, and set
$d_q(\mu,\mathcal A)=\min_{\nu\in\mathcal A}
[d_{\rm TV}(\mu,\nu)+W_q(\mu,\nu)]$ for finite nonempty $\mathcal A$.

\begin{corollary}\label{cor:distributional-stability}
Fix $p>1$, $p\ne2$, and $1\le q<\infty$. There is a finite,
nondecreasing function $\omega_{p,q}:[0,\infty)\to[0,\infty)$ with
$\omega_{p,q}(t)\to0$ as $t\downarrow0$ such that
\[
 d_q(\mu,\mathcal A_p)\le\omega_{p,q}(\Delta_p(\mu))
 \qquad(\mu\in\mathcal C).
\]
Moreover, if $K$ is a nonempty compact subset of one of
$(1,p^*)$, $(p^*,2)$ or $(2,\infty)$, a single function
$\omega_{K,q}$ has this property simultaneously for every $p\in K$.
\end{corollary}

\begin{proof}
We first record a uniform tail bound. For $X$ with law in $\mathcal C$,
let $S(t)=\Pr(X\ge t)$. This function is log-concave: apply preservation
under marginalisation to $f(t+u)\mathbf1_{\{u\ge0\}}$, where $f$ is
the density of $X$; see \cite[Theorem~3.3]{SW}.
Chebyshev's inequality gives $S(2)\le1/4$ and $S(-2)\ge3/4$.
If $S(2)>0$, the decreasing secant slopes of $\log S$ imply, for $t\ge2$,
\[
 S(t)\le S(2)\left(\frac{S(2)}{S(-2)}\right)^{(t-2)/4}
 \le\frac14\,3^{-(t-2)/4}.
\]
The bound is immediate if $S(2)=0$, and also holds where $S(t)=0$.
Applying it to $-X$ gives
\begin{equation}\label{eq:stability-uniform-tail}
 \sup_{\mu\in\mathcal C}\Pr_\mu(|X|\ge t)
 \le\frac12\,3^{-(t-2)/4}\qquad(t\ge2).
\end{equation}
In particular every fixed absolute moment is uniformly integrable on
$\mathcal C$.

The class is tight, so any sequence has a weakly convergent subsequence.
Log-concavity of probability measures is preserved under weak limits
\cite[Proposition~3.6]{SW}. By \eqref{eq:stability-uniform-tail}, the
limit retains mean zero and variance one. It is therefore nondegenerate
and has a log-concave density \cite[Theorem~2.7]{SW}, so lies in
$\mathcal C$. All fixed absolute moments converge. Convergence also
holds in $W_q$: couple by a common uniform quantile, which gives almost
sure convergence, and use the uniform higher moments to obtain
uniform integrability of the $q$th power of the coupling difference.
Finally, weak convergence to this nondegenerate log-concave density
implies convergence in total variation by
\cite[Proposition~2]{CS}. Its hypothesis that the limit has a density
has just been verified; no continuity of that density at its support
endpoints is required.

Suppose $\Delta_p(\mu_n)\to0$. Along any convergent subsequence as above,
moment continuity gives $\Delta_p(\mu)=0$. Theorem~\ref{thm:U} then
gives $\mu\in\mathcal A_p$, while convergence in both distances gives
$d_q(\mu_n,\mathcal A_p)\to0$ along this subsequence. A subsequence
bounded away from $\mathcal A_p$ is consequently impossible.
One may therefore take
\[
 \omega_{p,q}(t)=\sup\{d_q(\mu,\mathcal A_p):
                 \mu\in\mathcal C,\ \Delta_p(\mu)\le t\}.
\]
This function is finite by \eqref{eq:stability-uniform-tail}, and the
preceding argument proves its asserted limit (and $\omega_{p,q}(0)=0$).

For the uniform assertion, the set $\mathcal A_p$ is constant on $K$.
If uniformity failed, choose $p_n\in K$ and $\mu_n\in\mathcal C$
with $\Delta_{p_n}(\mu_n)\to0$ but distances bounded away from zero.
Pass to $p_n\to p\in K$ and a weakly convergent subsequence of $\mu_n$.
The tail bound implies joint continuity of $\E_{\mu_n}|X|^{p_n}$:
on bounded intervals the integrands converge uniformly, and the tails
are controlled by any fixed moment order larger than $\max K$.
The endpoint constants are continuous in $p$. Hence
$\Delta_p(\mu)=0$, and the same contradiction applies. Taking the
supremum also over $p\in K$ defines $\omega_{K,q}$.
\end{proof}

\paragraph{Limits of exponent-uniformity.}
The fixed-exponent statement includes $p=p^*$, with all three
variance-normalised laws in $\mathcal A_{p^*}$. It cannot, however,
be made uniform through $p^*$ with the exact exponent-dependent
equality sets: if $p_n\uparrow p^*$ with $p_n<p^*$, then
$\Delta_{p_n}(\nu_+)\to0$, but
$d_{\rm TV}(\nu_+,\nu_L)>0$ is fixed.
Nor is there such a uniform modulus as $p\to2$. For the standard
Gaussian law $\gamma$, $\Delta_p(\gamma)\to0$ as $p\to2$, whereas
$\min_{\nu\in\{\nu_+,\nu_-\}}d_{\rm TV}(\gamma,\nu)>0$.
At $p=2$ every law in $\mathcal C$ is an equality law.
Thus the corollary concerns convergence to the equality-law set, not
selection of a unique branch at the transition, and provides no
explicit rate in the moment deficit.

\section*{Supplementary material and reproducibility}
The bundle contains four supplementary notes with the complete closed forms, derivative identities, enclosure lemmas and box-level details used by the certificates: \texttt{supplement/bulk.md} (Propositions~\ref{prop:crossing}--\ref{prop:bulk}), \texttt{supplement/local.md} (Proposition~\ref{prop:local}), \texttt{supplement/strip.md} (Proposition~\ref{prop:strip}) and \texttt{supplement/largep.md} (Proposition~\ref{prop:largep}). It also contains the certifiers, their outputs, independent floating-point checks and a one-command replay script (\texttt{code/reproduce.sh}). All certificates use python-flint~0.9.0 (Arb) \cite{Johansson} with exact rational region boundaries, outward rounding, and failure on any undecided comparison.

\paragraph{Rigour of the library functions.} Arb's elementary and special functions used here (powers, $\exp$, $\log$, $\Gamma$, $\psi$, the upper incomplete gamma function, ${}_1F_1$, ${}_2F_2$ and Tricomi $U$) return, for ball inputs, a ball containing $f(x)$ for every $x$ in the input balls, or a non-finite ball when this cannot be guaranteed (for instance when an input ball meets a pole or a branch cut). Every certifier treats a non-finite output, like an undecided comparison, as a failure of the box. Integrals in the strip certificate use Arb's rigorous integration \cite{JohanssonInt}: the integrand callback returns enclosures on complex balls and signals non-analyticity (logarithms are evaluated with the analytic flag, so a ball touching the branch cut yields a non-finite value and the box fails). When an external parameter $s$ enters as a ball, every callback value encloses the integrand for all $s$ in that ball, so the resulting enclosure holds uniformly in $s$. Where monotonicity in a parameter is used for an exponent that may be negative (e.g.\ $\E(d+sR)^q$ with $q=p-2<0$ in L2), the certifier takes the hull of all corner evaluations, which is valid for either direction of monotonicity.

\paragraph{Proof-to-code map.}
The domains and exact sign statements are given in the cited propositions;
the following table identifies their executable counterparts. Paths in the
last two columns are relative to \texttt{code/} and \texttt{certificates/}.
All four supplements are included in the same public archive.
\begin{center}\small
\begin{tabular}{@{}lll@{}}
\toprule
Obligation & Certifier & Output \\
\midrule
Bulk B1--B4 (Prop.~\ref{prop:bulk}) & \texttt{bulk/bulk\_cert.py} & \texttt{bulk/cert\_B*.json}\\
Crossing (Prop.~\ref{prop:crossing}) & \texttt{bulk/crossing\_cert.py} & \texttt{bulk/cert\_X.json}\\
Local L1--L3 (Prop.~\ref{prop:local}) & \texttt{local/local\_cert.py} & \texttt{local/cert\_L*.json}\\
Strip S1--S2 (Prop.~\ref{prop:strip}) & \texttt{strip/strip\_cert.py} & \texttt{strip/S*\_cert.json}\\
Large $p$ (Prop.~\ref{prop:largep}) & \texttt{largep/largep\_certify.py} & \texttt{largep/largep\_cert.json}\\
\bottomrule
\end{tabular}
\end{center}
The large-$p$ conclusion is a written universal argument; the executable
checks its rational constants and optional diagnostics. Numerical samples do
not replace any universal claim. In particular the strip helper
\texttt{G\_direct} is a truncated diagnostic, not a certified enclosure of
the full function; production S1/S2 use the separately bounded routines.

\paragraph{Status and limitations.}
This is an unrefereed computer-assisted theorem candidate. Ball arithmetic
checks the stated regional inequalities, not the complete written argument
in a proof assistant. No unaffiliated reproduction or journal peer review is
asserted. The stability corollary is qualitative, and the general sharp
two-exponent moment problem remains open here. No arithmetic classification
of $p^*$, including irrationality, is claimed. The public archive is
identified by DOI \href{https://doi.org/10.5281/zenodo.23405173}{10.5281/zenodo.23405173}.

\begin{thebibliography}{99}
\bibitem{BP} S.~Brazitikos, C.~Pandis. Sharp inequalities for symmetric polynomials, Hunter's conjecture, and moments of exponential random variables. arXiv:2512.12254 (2025).
\bibitem{BTT} S.~Brazitikos, C.~Tang, T.~Tkocz. Moments of sums of exponentials, beyond CHS. arXiv:2602.03058 (2026).
\bibitem{Eitan} Y.~Eitan. The centered convex body whose marginals have the heaviest tails. \emph{Studia Math.}\ 274 (2024), no.~3, 201--215, doi:10.4064/sm211027-20-11; arXiv:2110.14382.
\bibitem{ENT} A.~Eskenazis, P.~Nayar, T.~Tkocz. Sharp comparison of moments and the log-concave moment problem. \emph{Adv.\ Math.}\ 334 (2018) 389--416.
\bibitem{Johansson} F.~Johansson. Arb: efficient arbitrary-precision midpoint-radius interval arithmetic. \emph{IEEE Trans.\ Comput.}\ 66 (2017) 1281--1292.
\bibitem{JohanssonInt} F.~Johansson. Numerical integration in arbitrary-precision ball arithmetic. In \emph{Mathematical Software -- ICMS 2018}, LNCS 10931, Springer (2018) 255--263.
\bibitem{MRTT} J.~Melbourne, M.~Roysdon, C.~Tang, T.~Tkocz. From simplex slicing to sharp reverse H\"older inequalities. \emph{J.\ London Math.\ Soc.}\ (2) 113 (2026), no.~3, e70461; arXiv:2505.00944v2.
\bibitem{PolyaSzego} G.~P\'olya, G.~Szeg\H{o}. \emph{Problems and Theorems in Analysis~II}. Springer, Berlin (1976).
\bibitem{Murawski} D.~Murawski. Comparing moments of real log-concave random variables. arXiv:2211.05210.
\bibitem{SW} A.~Saumard, J.~A.~Wellner.
Log-concavity and strong log-concavity: A review.
\emph{Statist. Surv.}\ 8 (2014), 45--114,
doi:\href{https://doi.org/10.1214/14-SS107}{10.1214/14-SS107}.
\bibitem{CS} M.~Cule, R.~Samworth.
Theoretical properties of the log-concave maximum likelihood estimator
of a multidimensional density.
\emph{Electron. J. Stat.}\ 4 (2010), 254--270,
doi:\href{https://doi.org/10.1214/09-EJS505}{10.1214/09-EJS505}.

\end{thebibliography}
\end{document}
